\hypertarget{appendix-aa405b7c3903}{}\label{referencestart:appendix-aa405b7c3903}
\section{Borrow contracts and pool epochs}\label{appendix-aa405b7c3903-borrow-contracts-and-pool-epochs}

\textbf{Curriculum depth:} Introduced · \textbf{First seen:} C008 · \textbf{Latest development:} C008

\subsection{Definition and mechanism}\label{appendix-aa405b7c3903-definition-and-mechanism}

An epoch versions a borrow-validity interval. A checked borrow compares its saved epoch to the live pool counter before using the cached element. The pool counter must outlive the borrow; counters alone do not supply ownership, synchronization, race freedom or wrap immunity.

\subsection{Prerequisites}\label{appendix-aa405b7c3903-prerequisites}

\begin{itemize}
\tightlist
\item
  \hyperref[appendix-f9b8b9188c7f]{Borrow invalidation and logical identity} --- Invalidating mutations define the deadlines a saved epoch or deferred borrow must respect.
\end{itemize}

\subsection{Used by}\label{appendix-aa405b7c3903-used-by}

Later sources will deepen this same node; no empty dependent nodes are created.

\subsection{Recall and next depth}\label{appendix-aa405b7c3903-recall-and-next-depth}

List every mutation that advances the epoch and identify check-then-use races. Pending: production wrappers, wrap policy and concurrent lifetime pinning. Written coverage does not establish learner mastery.

\subsection{Source and navigation}\label{appendix-aa405b7c3903-source-and-navigation}

\hyperref[01-foundations-i-ch003-when-a-valid-borrow-stops-naming-the-intended-thing]{Teaching explanation and examples} · \hyperref[appendix-fd090d982a15-m008---parts-113-126]{Exact coverage} · \hyperref[appendix-4ad47bdd269a]{Dependency Map}
